\section{Introduction}

% Discuss rise of LLMs in automated system config/API synthesis
% Shift from HitL to agentic workflows/AI agents (ie; x402)
\subsection{Motivation}
In recent years, the deployment of Large Language Models (LLMs) has transitioned
from simple, explanatory or otherwise "conversational" interfaces, to highly
mission-critical integrations with autonomous AI and agentic workflows~\cite{exampleRef}.
This shift towards "AI agents" and automated system configuration---exemplified
by frameworks such as AutoGPT, OpenClaw, and the x402 protocol---is indicative
of an increasing reliance on these models to synthesize complex API calls and
cloud infrastructure definitions with minimal human-in-the-loop verification
~\cite{exampleRef}. While the transformer architecture~\cite{exampleRef} has demonstrated
remarkable fluency in natural language, its application in safety-critical
structural synthesis remains somewhat dubious. In industry, the move towards
fully autonomous agents requires that the output not only "looks" like valid
code, but is also actually \textit{compliant} with the operational state of the
target system---that is to say, beyond just syntax.

% Boundary between CFG/CSI (context-sensitive invariants)
% Why does FSM-based decoding fail? State dependent logic needed?
\subsection{Problem Statement}
Current methodologies for steering LLM output primarily rely on constrained
decoding via Finite State Machines (FSMs) or Context-Free Grammars
(CFGs)~\cite{exampleRef}. While these techniques effectively \textit{guarantee}
surface-level syntax (e.g., ensuring a bracket follows a key in a JSON object),
they are fundamentally incapable of enforcing Context-Sensitive Invariants
(CSIs).

Naturally, there is an important distinction between pure syntactic validity and
semantic correctness that modern approaches do not capture. For instance, an FSM
can ensure a configuration file is "well-formed" JSON, but it cannot verify that
an assigned IP address is within a previously declared subnet range, or that a
\texttt{start\_date} chronologically precedes an \texttt{end\_date}. These are
state-dependent constraints aligned with Type-1 Chomsky hierarchy
machines~\cite{exampleRef}, and are therefore beyond the reach of context-free parsers
utilized by modern tools such as Outlines or Guidance~\cite{exampleRef}. As a consequence,
when LLMs encounter these "semantic walls", they either hallucinate invalid
state transitions, leading to system-wide failures during execution, or post-hoc
checks demand another attempt, significantly increasing system latency and
computation cost.

% Introduce dual-process architecture
\subsection{Proposed Contributions}
Thus, this research proposes a dual-process neuro-symbolic framework designed to
provide formal structural guarantees while maintaining a reasonable inference
efficiency. The crux of this work is the Syntax Oracle, an abstract verification
machine that maintains the specification state during the decoding process. By
projecting invalid neural proposals onto an admissible constraint space via
dynamic logit masking, the Oracle acts as a deterministic safety controller.

Simultaneously, we introduce a Semantic Internalization Pipeline, where rather
than relying solely on the Oracle to correct runtime errors, we can utilize it
as a "teacher" to generate synthetic traces used to train a model for
prediction. That is to say, we use these traces to fine-tune the LLM using
Low-Rank Adaptation (LoRA)~\cite{exampleRef}. 

Additionally, we can quantify the efficacy of this system by evaluating the
Oracle Intervention Rate (OIR)---a metric to measure the model's learned
structural competence and ability to reduce reliance on the deterministic
failsafe oracle by looking at the rate at which the oracle must employ its
masking. Naturally, other industry-standard benchmarks will also be used.

